\section{Persistent marks and guarded history routes}
\label{ba:marked-routes}

The preceding recorders clear the work-head mark at each checkpoint.
We now retain that mark, or replace the moving history pointer by a
guarded boundary between occupied and unused cells.  These changes give
different partial maps on arbitrary tapes, even though their initialized
executions simulate the same machine.  We prove each inverse on its
declared domain before using it to separate the geometric branch images.
The use of retained history follows the reversible-computation idea of
Bennett~\cite{Bennett1973}; the local guards and geometric interfaces
are established here.

In this section an input word occupies cells $0,1,\ldots$, the other
work cells are blank, and the initial head is at zero.  All constructions
include initial halting.  When one halt state $h$ is used, identify the
original halting states with $h$ and redirect incoming instructions;
there are no outgoing halting instructions to preserve.  If necessary
adjoin an unreachable halt state.
We write $X(t;a)=\Phi_t(a)$ for the trajectory from $a$.

\subsection{A persistent mark and input-dependent placement}
\label{ba:persistent-periodic}

A force can be periodic from time zero if no separate loading interval
is needed.  We achieve this by translating the initial control's coding
square so that the prescribed particle already represents the input.
The translation depends effectively on the finite input.  First we give
the recorder that this periodic processor implements.

\begin{lemma}[A recorder with a persistent work-head mark]
\label{ba:persistent-compiler}
Let a machine have work alphabet $\Gamma$, halt state $h$, and
instructions
$\gamma=(q,a)\mapsto(q^+_\gamma,a^+_\gamma,d_\gamma)$,
where $q\ne h$ and $d_\gamma\in\{-1,0,1\}$.
There is an effective finite partial one-head table with these properties.
At the checkpoint after $n$ work steps, it has the correct work tape,
state and head, a unique mark at that head, and a history pointer at
$p_n=n+2$.  Each target control determines the incoming displacement;
the target control and written full cell symbol determine the source
control and read full symbol on the entire partial domain.  If the next
work-head position is $w'=w_n+d_\gamma$, the next checkpoint is reached
after exactly $2(p_n-w')+4$ local steps.
\end{lemma}

\begin{proof}\label{ba:persistent-compiler-proof}
A full symbol is $(a,e,s)\in\Gamma\times\{0,1\}\times\mathcal H$,
where $\mathcal H=\{E,P\}\sqcup\{r_\gamma\}$ records empty history,
a pointer, or an instruction.  All displayed control families are
disjoint.  The following is the complete table; free cell indices range
over their finite alphabets:
\[
\begin{array}{c|c|c|c|r}
\text{control}&\text{read and guard}&\text{target}&\text{write}&d\\\hline
A_q&(a,1,s),\ s\ne P&B_\gamma&(a^+_\gamma,0,s)&d_\gamma\\
B_\gamma&(c,0,s),\ s\ne P&R_\gamma&(c,1,s)&1\\
R_\gamma&(c,0,s),\ s\ne P&R_\gamma&(c,0,s)&1\\
R_\gamma&(c,0,P)&C_{q^+_\gamma}&(c,0,r_\gamma)&1\\
C_q&(c,0,E)&L_q&(c,0,P)&-1\\
L_q&(c,0,s),\ s\ne P&L_q&(c,0,s)&-1\\
L_q&(c,1,s),\ s\ne P&A_q&(c,1,s)&0
\end{array}
\]
The first row has $q\ne h$ and $\gamma=(q,a)$; there is no rule from
$A_h$.  Initially the actual head and the unique mark are at zero,
the work tape is the input, and history is $E$ except for $P$ at cell 2.

At checkpoint $n$, $|w_n|\le n$, so after the first row the head is at
$w'\le n+1<p_n$.  That row clears the old mark, and the second installs
the new one at $w'$.  The right scan reaches the pointer without meeting
another mark, writes $r_\gamma$, and visits the empty cell $p_n+1$.
The next row installs the new pointer there.  The left scan returns
through the newly written record to the unique mark and enters
$A_{q^+_\gamma}$ without clearing it.  Hence every invoked rule is
defined, the work instruction is correct, and the checkpoint invariant
is restored.  The two continuing scans take $p_n-w'-1$ and $p_n-w'$
steps; the other five steps give the stated total.  This also handles
$d_\gamma=0$.  Induction gives a defined execution until and including
the first halting checkpoint, and indefinitely in the nonhalting case.

We check the inverse independently of this initialization.  Targets
$B_\gamma,R_\gamma,C_q,L_q,A_q$ have incoming displacements
$d_\gamma,1,1,-1,0$, respectively.  Undo that displacement to locate
the written cell.  At $B_\gamma$, the index $\gamma$ recovers the old
control and overwritten work symbol, and the old mark was one.  At
$R_\gamma$, the written mark is one for entry and zero for the scan,
so it identifies the source and the old mark.  At $C_q$, the written
record identifies $\gamma$ and its former value $P$.  At $L_q$, a
written pointer identifies entry, whereas every loop has history
different from $P$.  At $A_q$, the unique source control is $L_q$ and
the full symbol is unchanged.  All unspecified tracks are retained.
Thus every possible output has at most one preimage on the full domain.
\end{proof}

\begin{theorem}[An input-dependent entry chart]
\label{ba:periodic-entry}
For each machine and finite input, and each fixed positive computable
viscosity $\nu$, there is an effective smooth force on $\mathbb R^3$
with one compact spatial support, bounded mixed derivatives, and period
one from time zero, whose zero-data solution satisfies
\[
 \exists t\ge0:\ X_1(t;(2,0,0))<0
 \quad\Longleftrightarrow\quad\text{the machine halts}.
\]
The solution is smooth with pressure zero and is unique among smooth
solutions for which, on every finite interval,
$u\in C_tH^2\cap C_t^1L^2$, $u$ and $\nabla u$ are bounded, and
$p\in C_tH^1$.
As a separate force choice, the same event is realized with every
mixed derivative of velocity and force in space-time $L^2$ and bounded
by a derivative-dependent multiple of $(1+t)^{-1}$.
\end{theorem}

\begin{proof}\label{ba:periodic-entry-proof}
Prepend a nonhalting state $q_{\rm in}$ which, on every work symbol,
preserves it, stays put, and enters the original initial state.  This
includes original initial halting while keeping the new initial state
nonhalting.  Apply Lemma~\ref{ba:persistent-compiler}.  For its full
alphabet of size $m$ use odd digits $e(\sigma)\in\{1,3,\ldots,2m-1\}$
in base $B=2m+2$, and write
\[
 c(\omega)=\sum_{j\ge0}e(\omega_j)B^{-j-1},\qquad
 (\xi,\eta)=
 \big(c(\sigma_{-1}\sigma_{-2}\cdots),
       c(\sigma_0\sigma_1\cdots)\big).
\]
The initial values $\xi_0,\eta_0$ are rational, because the initialized
streams have finite prefixes and known constant tails.  Set
\[
 O(A_{q_{\rm in}})=(2-\xi_0,-\eta_0),\qquad
 O(A_h)=(-2,0),\qquad O(P)=(4+2j,0)
\]
for the remaining controls in a finite ordering.  All squares
$O(P)+[0,1]^2$ are separated.  Every nonhalting square lies strictly
to the right of $x_1=1$, and the halting square is
$[-2,-1]\times[0,1]$.  The initial code at height zero is exactly
$(2,0,0)$.

Apply Lemma~\ref{lem:bb-onecell} with digits $d(\sigma)=e(\sigma)$,
base $B=2m+2$, and the state translations just specified.  Its incoming
hypotheses are exactly those of Lemma~\ref{ba:persistent-compiler}.
With $I_\sigma=(e(\sigma)+[0,1])/B$, use the minimal left-split
branches of \eqref{eq:bb-one-cell-branches}, including every left
letter for a left move.  The state translations conjugate these maps
without changing their reciprocal diagonal linear parts.  The shared
lemma therefore gives the exact updates and separate source and target
separation on full closed rectangles.  If an absolute head position is wanted at a checkpoint,
the number of contiguous records gives $n$, and the pointer's relative
position locates the head relative to its known absolute position $n+2$.
The gapped digits permit each finite prefix needed for this recovery
to be read with increasing precision.

Apply the exponential interpolation of
Lemma~\ref{ba:two-interpolations}, with heights $10i$ and the following
separated phase intervals: lifting on $[1/8,2/8]$, horizontal motion
on $[3/8,5/8]$, and descent on $[6/8,7/8]$.  On each interval use a
rescaled flat step $\sigma$ from \eqref{eq:p2-step}.
Let $\delta$ be one quarter of the minimum of 1
and the positive planar source and target gaps.  Use plateau padding
$\delta/3$ and support padding $2\delta/3$.  The same gap and height
estimates in the lemma make these collars disjoint.  Here the local
potential can be written
\[
 \tfrac12 C_i'\times(x-C_i)
 +(0,0,\eta'\log\lambda_i\,(x_1-C_{i1})(x_2-C_{i2})).
\]
Its curl is the required translation plus
$\eta'\log\lambda_i\,(x_1-C_{i1},-(x_2-C_{i2}),0)$.
It realizes the whole cuboid and a neighborhood by the plateau argument.

Repeat this unit-time field as $U$, and put
$f=U_t+(U\cdot\nabla)U-\nu\Delta U$ at the prescribed physical
viscosity.  The zero collars give $U(0)=0$ and smooth period-one gluing.
Lemma~\ref{ba:realization} gives the solution, uniqueness, support and
effective mixed-derivative bounds.  Its hypotheses hold because this is
a finite curl formula on a compact spatial set and one compact time
period.  The formula includes every branch and never requires the run
on the selected input.

At integer times, induction gives the exact local configuration through
the first halt, or at every integer time for a nonhalting run.  A halt
therefore reaches the negative square.  During a nonhalting run, both
endpoint squares of every used branch are nonhalting.  Every moving
center has first coordinate greater than one and every first half-width
is at most $1/2$, because the endpoint widths are at most one and the
exponential scale is monotone between them.  The actual particle stays
positive during all phases.  This proves the equivalence for all real
times, including the added entry step.

For the separate decaying choice apply Lemma~\ref{ba:log-clock}:
$\widehat U(t,x)=(1+t)^{-1}U(\log(1+t),x)$, with its own residual at
the same $\nu$.  That lemma gives the stronger temporal-order-$k$
bound $O((1+t)^{-1-k})$ for every spatial derivative, hence the stated
bound and space-time $L^2$ conclusion on the common compact support.
The clock is onto $[0,\infty)$, and the chain rule identifies the new
trajectory with $X(\log(1+t);(2,0,0))$.  The fixed event is unchanged.
\end{proof}


\subsection{Using a fresh half-line instead of a frontier symbol}

The next rule reads a neighboring history cell.  Its inverse therefore
uses a full three-cell cylinder, rather than only the written cell.

\begin{lemma}[Three-cell fresh-tail recorder]\label{ba:fresh-recorder}
There is an effective partial rule on cells at offsets $-1,0,1$ that
simulates the machine, is injective on its entire domain of relative
tapes, and has incoming displacement determined by the target state.
Its allowed cylinders admit separately separated closed-rectangle maps
with linear parts $\operatorname{diag}(B^{-d},B^d)$.
\end{lemma}
\begin{proof}\label{ba:fresh-recorder-proof}
Use symbols $c_i=(a_i,\ell_i,\mu_i)$ in
$C=A\times(\{\star,\perp\}\sqcup\mathcal R)\times\{0,1\}$
and controls $N_q,B_q,J_r,S_r$.  Only $N_q$ with $q\in H$ halt.
After writing, a move $d$ reindexes a relative tape by
$c'_j=\bar c_{j+d}$.  The following rules give the complete domain:
\begin{enumerate}
\item At $N_q$, $q\notin H$, require $\mu_0=0$ and
$\ell_0\ne\perp$.  With $r=(q,a_0)$ write $b(r)$, move $d(r)$,
and enter $J_r$.
\item At $J_r$, require $\mu_0=0$ and $\ell_0\ne\perp$.
Mark cell zero, move right, and enter $S_r$.
\item At $S_r$, require $\mu_0=0$.  If $\ell_0\ne\perp$, scan
right.  If $\ell_0=\ell_1=\perp$, write record $r$ at zero,
move left, and enter $B_{q^+(r)}$.
\item At $B_q$, require $\ell_0,\ell_1\ne\perp$.
Scan left if $\mu_0=0$; if $\mu_0=1$, clear the mark, stay, and
enter $N_q$.
\end{enumerate}
Initialize history as $\star$ below absolute cell 2 and $\perp$ from
2 onward, all marks zero, the prescribed work tape, and control $N_{q_0}$.
At checkpoint $n$ the first fresh cell is $F=n+2$, records occupy
$2,\ldots,F-1$, and the work head satisfies $h\le F-2$.
The work move reaches $k\le F-1$.  Mark $k$, scan to $F$, write the
record, and return from $F-1$ to $k$.  The return's two-history guard
holds even at its first cell because cell $F$ has just been filled.
Clearing the mark restores the invariant with frontier $F+1$.
There are $F-k-1$ instructions in each continuing scan and four others,
so the duration is $2(F-k)+2$.  This proves finite completion and the
checkpoint halt equivalence.

For the full-domain inverse, targets $J_r,S_r,B_q,N_q$ determine moves
$d(r),1,-1,0$, respectively.  Undo the indicated shift.  At $J_r$,
$r$ recovers the overwritten letter and old control.  At $S_r$, the
written mark distinguishes entry from the zero-mark scan.  At $B_q$,
target offset $+2$ is the former offset $+1$; its history is fresh
for a recording entry and nonfresh for a return scan.  In the entry
case target offset $+1$ is the former scanned cell; its record identifies $r$.
At $N_q$, restore the mark to one.  All other tracks are preserved.
The neighbor guard is what makes this inverse valid on arbitrary allowed
tapes.

Use odd digits in base $B=2|C|+1$ and state offsets $e_s$ whose
unit intervals have gaps at least two.  A permitted tuple
$(s,c_{-1},c_0,c_1)$ has source
$(e_s+I_{c_{-1}})\times I_{c_0c_1}$.  The full inverse just proved
checks the injectivity hypothesis of Lemma~\ref{lem:bb-window-rectangles}
with $(a,b)=(1,2)$ and $d\in\{-1,0,1\}$.
Only cell zero is written, which is within that window.  The lemma
therefore gives full target cylinders with prefix lengths
$(1+d,2-d)$, exact affine factors $(B^{-d},B^d)$, and separate
closed-rectangle images.  For one target state the incoming $d$ is
fixed, so its prefix lengths agree across incoming branches.  An
incompatible digit among at most three positions yields the additional
quantitative target gap $B^{-3}$.  This is separation of the filled
rectangles, including their boundaries.
\end{proof}

\begin{theorem}[Compact realization of fresh-tail cylinders]
\label{ba:fresh-flow}
The fresh-tail recorder has an effective compactly supported smooth
realization on $\mathbb R^3$, with zero initial velocity and force
periodic after one unit of loading, whose event
$\exists t\ge0:(\Phi_t(0))_1<-1$ is equivalent to halting.
A separate choice, with the same event and one compact support, satisfies
for every $k\ge0$ and spatial multi-index $\alpha$
\begin{equation}\label{ba:fresh-decay}
 \sup_x\bigl(|\partial_t^k\partial_x^\alpha u|+
             |\partial_t^k\partial_x^\alpha F|\bigr)
 \le C_{k,\alpha}(1+t)^{-1-k}.
\end{equation}
Every displayed derivative is space-time square-integrable.  Both forces
have the unique zero-pressure solution of Lemma~\ref{ba:realization}
at the fixed computable viscosity.
\end{theorem}
\begin{proof}\label{ba:fresh-flow-proof}
Use the full alphabet $C$ of Lemma~\ref{ba:fresh-recorder} and
$B=2|C|+1$.  Place halting controls at $e_s=-3j_s$ and all others
at $3j_s$, with distinct positive integer labels.  Use the three-cell rectangles as
plates at height zero.  Branch $i$ has middle height $4i$ and padding
$B^{-3}/10$.  Interpolate both planar factors independently and apply
Lemma~\ref{ba:plate-routing}, including its reciprocal normal strain.
Planar gaps separate the vertical phases; middle heights separate the
padded plates, whose collar half-thickness is less than one.  Thus the
specified motion holds on full neighborhoods.  It retains independent
planar interpolation throughout the cycle.

The initial code $P_*$ is rational.  A point-centered curl translation,
with center $\theta(t)P_*$ and padding one, loads zero in one unit.
Repeat the plate cycle thereafter.  The finite-step invariant and code
maps identify the trajectory at all times $1+n$ until halt and for all
$n$ otherwise.  A halt code has first coordinate below $-1$.  For a
nonhalting run, loading is nonnegative and each later pair of coded
endpoints is positive; the two-coordinate convex interpolation of
Lemma~\ref{ba:plate-routing} proves the all-time exclusion.
Use residual realization for the periodic choice and
Lemma~\ref{ba:log-clock} for the separate decay choice.
\end{proof}


\subsection{Guarded history cylinders in a torus chart}
\label{ba:guarded-history}

An occupied prefix can replace the explicit pointer.  The resulting
inverse must distinguish arrival from the recording cell from an
ordinary return-scan step.  Testing the next history cell supplies that
distinction.  Unlike the fresh-tail recorder, the marked return below
does not test its right neighbor; we retain this larger partial domain.

\begin{lemma}[An occupied-prefix history recorder]
\label{ba:occupied-compiler}
Let a machine have state set $Q$, halting set $H$, alphabet $\Gamma$,
and instructions
$\tau=(q,a)\mapsto(q^+_\tau,w_\tau,d_\tau)$.
There is an effective partial injective map on head-relative tapes
whose domain depends only on the control and full cells at offsets
$-1,0,1$.  With the initialization below, occupied history at checkpoint
$n$ consists exactly of the absolute indices less than $F_n=n+2$.
A work step ending at head position $j$ takes exactly
$2(F_n-j)+2$ local steps.  Halting at a checkpoint, including the initial
one, is equivalent to halting of the original machine.
\end{lemma}

\begin{proof}\label{ba:occupied-compiler-proof}
Full symbols are $(a,h,m)$, with work letter $a$, mark $m\in\{0,1\}$,
and history $h\in\{\#,E\}\sqcup\mathcal R$, where
$\mathcal R=(Q\setminus H)\times\Gamma$.  History is called occupied
when $h\ne E$.  A move $d$ reindexes the post-write relative tape as
$\xi^{\rm new}_k=\widetilde\xi_{k+d}$.  The complete table is
\[
\begin{array}{c|l|l|c|r}
\text{control}&\text{guard}&\text{write at }0&\text{target}&d\\\hline
\mathsf{Run}_q&m_0=0,\ h_0\ne E&a_0\leftarrow w_\tau&
 \mathsf{Mark}_\tau&d_\tau\\
\mathsf{Mark}_\tau&m_0=0,\ h_0\ne E&m_0\leftarrow1&
 \mathsf{Seek}_\tau&1\\
\mathsf{Seek}_\tau&m_0=0,\ h_0\ne E&\text{unchanged}&
 \mathsf{Seek}_\tau&1\\
\mathsf{Seek}_\tau&m_0=0,\ h_0=h_1=E&h_0\leftarrow\tau&
 \mathsf{Back}_{q^+_\tau}&-1\\
\mathsf{Back}_q&m_0=0,\ h_0\ne E,\ h_1\ne E&\text{unchanged}&
 \mathsf{Back}_q&-1\\
\mathsf{Back}_q&m_0=1,\ h_0\ne E&m_0\leftarrow0&
 \mathsf{Run}_q&0.
\end{array}
\]
The first row requires $q\notin H$ and $\tau=(q,a_0)$; all unmentioned
tracks are unchanged.  Initialize the original work tape with head zero,
all marks zero, history $\#$ at every index below 2, and $E$ elsewhere.

At checkpoint $n$ the original head has $|j_n|\le n$.  The work step
ends at $j\le n+1<F_n$, on occupied history.  Mark that cell and scan
right across occupied unmarked cells to $F_n$.  Both $F_n$ and $F_n+1$
are unused, so the recording row applies and moves left.  Every unmarked
return-loop cell has occupied history at itself and its right neighbor,
including the new record at $F_n$.  The return reaches $j$, clears the
mark, and enters the new work state.  The continuing scans each take
$F_n-1-j$ steps, with four other steps.  This proves finite completion,
the duration, and the new occupied boundary $F_n+1$.  Induction gives
the stated computation and halt equivalence.

For the full-domain inverse, target controls determine displacement.
The moves are $d_\tau$ into $\mathsf{Mark}_\tau$, one into
$\mathsf{Seek}_\tau$, minus one into $\mathsf{Back}_q$, and zero
into $\mathsf{Run}_q$.  At $\mathsf{Mark}_\tau$, undoing the
move and restoring the read work symbol encoded by $\tau$ recovers the
predecessor.  At $\mathsf{Seek}_\tau$, target offset $-1$ is the
written cell; its mark distinguishes the marking entry from the scan.
At $\mathsf{Back}_q$, target offset 2 is the former offset 1.  Its
history is $E$ for a recording entry and is occupied for a return loop.
In the recording case, target offset 1 contains $\tau$, identifying
the old control and restoring its history to $E$.  Finally, the only
entry into $\mathsf{Run}_q$ clears a mark without moving, so the mark
is restored to one.  These inverses use only the table's guards and
retained symbols, not the special shape of initialized history.
\end{proof}

\begin{lemma}[The full cylinders of the guarded recorder]
\label{ba:occupied-cylinders}
The recorder of Lemma~\ref{ba:occupied-compiler} admits a finite
effective list of positive-area closed source and target rectangles,
separately pairwise disjoint.  If $\mathcal A$ is its full alphabet and
$B=2|\mathcal A|+1$, a branch of displacement $d$ has an exact affine
map with linear part $\operatorname{diag}(B^{-d},B^d)$.
\end{lemma}

\begin{proof}\label{ba:occupied-cylinders-proof}
Apply Lemma~\ref{lem:bb-window-rectangles} with window parameters
$(a,b)=(1,2)$ to the complete table of
Lemma~\ref{ba:occupied-compiler}.  That table reads only offsets
$-1,0,1$, writes only offset zero, and has displacements
$d\in\{-1,0,1\}\subset[-1,2]$; its full-domain injectivity was
proved above.  Thus every hypothesis of the window lemma holds.

To fix the actual rectangles, choose odd digits
$\gamma(a)\in\{1,3,\ldots,2|\mathcal A|-1\}$ and
$B=2|\mathcal A|+1$.  Write
\[
 e(P_1\cdots P_k)=\sum_{j=1}^k\gamma(P_j)B^{-j},\qquad
 I(P)=e(P)+[0,B^{-k}],\qquad I(\varnothing)=[0,1],
\]
using the infinite series for an infinite word.  Choose separated
state intervals $o_s+[0,1]$ and use the code
\[
 r(s,\xi)=
 (o_s+e(\xi_{-1}\xi_{-2}\cdots),\ e(\xi_0\xi_1\cdots),0).
\]
For each permitted triple $(l,a,b)$ at offsets $-1,0,1$, let $a'$
be its written full symbol.  The source prefix pair is $((l),(a,b))$,
and the target pair is
\[
 \begin{cases}
 ((a',l),(b)),&d=1,\\
 ((l),(a',b)),&d=0,\\
 (\varnothing,(l,a',b)),&d=-1.
 \end{cases}
\]
Take the corresponding products of prefix intervals with the source
and target state offsets.  Both arbitrary tails are unchanged, so the
window lemma gives the exact center-to-center affine map with factors
$B^{-d},B^d$, onto the whole specified target cylinder and filled
rectangle.  Its separation conclusion applies to these complete closed
rectangles, including their boundaries.  All endpoints and positive
within-family gaps are effectively rational.  The odd digits also put
all infinite-word codes strictly in $(0,1)$ in each normalized coordinate.
\end{proof}

\begin{theorem}[A guarded processor in a fixed torus chart]
\label{ba:occupied-torus}
For each machine and finite input, at any fixed positive computable
viscosity there is an effective smooth mean-zero force on the unit flat
three-torus, with bounded mixed derivatives and period one for $t\ge1$,
whose smooth zero-data solution has pressure zero and satisfies
\[
 \exists t\ge0:\ X(t;[(-1/16,0,0)])\in
 \{[x]:x_1\bmod1\in(0,1/2)\}
 \quad\Longleftrightarrow\quad\text{the machine halts}.
\]
The pressure is normalized to mean zero, and the solution is unique in
the classical torus comparison class of Section~\ref{sec:model}.
\end{theorem}

\begin{proof}\label{ba:occupied-torus-proof}
Use the occupied-prefix recorder.  Assign offsets $2j$, $j\ge1$, to
halting run controls and $-2j$ to all other controls, with distinct
indices within each group.  The initial code $r_{\rm in}$ is rational:
the left stream is constant and the right stream is eventually constant.
If $N$ is the number of cylinders, set
\[
 O=1+\max_s|o_s|,\qquad K=O+6\max(1,N)+4,\qquad
 \sigma=(8K)^{-1},\qquad b_*=(-K/2,0,0).
\]
Use the linear version of Lemma~\ref{ba:two-interpolations} for the
branch boxes of third interval $[-1,1]$, obtaining $V_1$.
Use the same lemma with scale one on the single pair of cubes of
half-width one centered at $b_*$ and $r_{\rm in}$, obtaining a loader
$V_0$.  Extend both fields by zero outside their unit time intervals.
Their supports lie in $(-K,K)^3$.  For $V_1$, first-coordinate centers
have magnitude at most $O$, second-coordinate centers lie in $[0,1]$,
horizontal half-widths are at most $1/2$, and heights are at most $6N$.
For $V_0$, the horizontal center is bounded by $\max(K/2,O)$,
half-widths are one, and height is at most six.  The padding is at most
$1/4$ in either application.  Each bound is strictly less than $K$.

Periodize the rescaled fields in space and the step field in time:
\[
 U(t,[x])=\sigma\sum_{k\in\mathbb Z^3}
 \left[V_0(t,(x-k)/\sigma)+
       \sum_{n\ge1}V_1(t-n,(x-k)/\sigma)\right].
\]
Spatial support lies in chart boxes of radius $1/8$, so distinct spatial
copies are disjoint.  The time collars make the gluing smooth, initially
zero and periodic after time one.  Every mixed derivative is bounded.
The rescaled blocks remain curls of compactly supported potentials,
so $U$ has zero spatial mean.  Define the force at the physical viscosity
by $f=U_t+(U\cdot\nabla)U-\nu\Delta U$.  It has zero mean because
the convection term is $\operatorname{div}(U\otimes U)$ and the
Laplacian integrates to zero.  Lemma~\ref{ba:realization} now gives
the unique solution $(U,0)$.  Recomputing the residual after the chart
scaling is essential: it does not identify two different viscosities.

The prescription is effective without deciding a space or time seam.
Given approximations $\widehat x,\widehat t$ with error at most $1/4$,
all possibly active spatial indices satisfy
$\|\widehat x-k\|_\infty\le1$, and repeated time indices satisfy
$|\widehat t-n|\le2$.  Evaluate this finite superset using the flat
cutoffs and their derivative bounds.  The finite table, not an executed
machine history, determines every coefficient.

Since $\sigma b_*=(-1/16,0,0)$, loading sends the fixed particle to
$[\sigma r_{\rm in}]$.  At time $1+j$ it is the code after $j$ local
steps, through the first halt or forever in the nonhalting case, by the
exact whole-box maps.  Halting, including initial halting, gives a
positive first-coordinate representative below $1/8$, inside the target
half-circle.  For a nonhalting input both loading endpoints and both
endpoints of every subsequent used branch have negative first coordinate.
The linear-scale identity keeps the first coordinate negative at every
intermediate time.  All representatives stay in $(-1/8,0)$, disjoint
from the target modulo one.  This proves the all-time equivalence.
\end{proof}
